\section{Welfare and Coordination}
\label{sec:welfare}

The main contribution concerns strategic instrument switching, not a generic decentralization inefficiency. Nevertheless, the model delivers a useful welfare benchmark because the decentralized governments internalize different margins from a coordinated policymaker.

Let the coordinated policy vector maximize
\begin{equation}
\max_{s_A,h_A,s_B,h_B}\;W_A+W_B.
\label{eq:planner}
\end{equation}
For any policy $z_A\in\{s_A,h_A\}$, the difference between the coordinated and decentralized first-order conditions can be written as
\begin{equation}
\Omega_{z_A}^{W}
\equiv
\frac{\partial W_B}{\partial z_A}.
\label{eq:welfare-wedge-def}
\end{equation}
Using the structure of $W_B$,
\begin{equation}
\Omega_{z_A}^{W}
=
\frac12\frac{\partial CS}{\partial z_A}
+
\frac{\partial PS_B}{\partial z_A}
-
d(E_B-\bar E_B)\frac{\partial E_B}{\partial z_A}.
\label{eq:welfare-wedge}
\end{equation}
The wedge therefore contains three real cross-jurisdictional effects: consumer-surplus effects, producer-surplus/business-stealing effects, and emissions-target effects. The welfare comparison is not generated by subsidy transfers alone.

The sign of each wedge is parameter dependent, so the theory does not imply a global ranking of each decentralized policy instrument relative to its coordinated counterpart. The canonical interior witness nevertheless illustrates a systematic portfolio difference. Table~\ref{tab:welfare} reports symmetric decentralized and coordinated policies for the parameter vector in Section~\ref{sec:results}.

\begin{table}[htbp]
\centering
\caption{Decentralized and coordinated policies in the canonical witness}
\label{tab:welfare}
\begin{tabular}{cccccc}
\toprule
$\theta$ & $s^{NE}$ & $h^{NE}$ & $h^{NE}/(s^{NE}+h^{NE})$ & $s^{C}$ & $h^{C}/(s^{C}+h^{C})$\\
\midrule
0.60 & 2.941 & 1.143 & 0.280 & 7.003 & 0.396\\
0.774 & 2.444 & 0.941 & 0.278 & 5.406 & 0.406\\
0.90 & 2.059 & 0.744 & 0.265 & 4.525 & 0.414\\
1.00 & 1.709 & 0.524 & 0.235 & 3.923 & 0.422\\
\bottomrule
\end{tabular}
\begin{flushleft}\footnotesize
Notes: $h^C$ equals 4.597, 3.691, 3.199, and 2.867 at the four rows, respectively. The table is an illustration, not a global policy-ordering theorem.
\end{flushleft}
\end{table}

At $\theta=0.9$, total welfare is approximately $2.187$ in the decentralized equilibrium and $3.720$ under coordination. The infrastructure share rises from $0.265$ to $0.414$. Similar portfolio tilts appear across the displayed rivalry levels. The economic interpretation is that a coordinated policymaker internalizes the cross-jurisdictional consequences of the two policy instruments, whereas each local government chooses its portfolio to protect local welfare given the rival's policy.

These numbers should not be interpreted as a recommendation for globally maximizing carbon policy. The model's emissions term represents territorial or delegated decarbonization objectives. A global climate externality depending on $E_A+E_B$ would introduce an additional coordination motive. We deliberately leave that global-damage extension outside the baseline so that the paper isolates the strategic interaction between local green industrial policy and product-market rivalry.
